\unnumberedChapter{Introduction}

Designers of CPUs in the seventies and early eighties almost all used a microcoded approach with hundreds of instructions. RISC was just beginning to emerge. Registers were typically not fully general-purpose but often had special functions or meanings, and many instructions were rather complicated, though designers believed them useful. Compiler writers, however, often used only a small subset, ignoring these complex instructions because they were too specialized. In the late eighties and nineties the shift moved to RISC-based designs, with large sets of registers, fixed instruction word lengths, large virtual memory address ranges, and instructions that were in general much more pipeline-friendly. Today, processors are highly complex, with superscalar deep pipelines and multilevel cache hierarchies—just a few of the hallmarks of modern CPU design. CPU designs have become so large and complex that no single person can completely understand them. But what if there were a simple design that one person could fully understand?

\unnumberedSection{Twin-64}

Welcome to Twin-64. Twin-64 is a simple 64-bit CPU with a register–memory model and segmented virtual memory. The design is heavily influenced by Hewlett-Packard’s PA\_RISC architecture, which was initially a 32-bit RISC-style load/store machine. Many of the key architectural concepts originated there. Other processors, including DEC Alpha, IBM PowerPC, and more recently RISC-V, also influenced the design or were studied for their architectural approaches. In addition, the Blitz-64 architecture, developed at Portland State University, provided inspiration through its exclusive use of 64-bit signed arithmetic, thereby avoiding many of the problems associated with mixing signed and unsigned operations.

Three design principles guided the development of the Twin-64 architecture. First, the instruction set should be as regular and orthogonal as possible, making instructions easy to understand, decode, and combine. Second, the hardware implementation should remain simple and predictable, avoiding unnecessary complexity where practical. Third, instruction-level parallelism should be exposed primarily by software rather than discovered dynamically by hardware. These principles influenced both the instruction set architecture and the processor implementations presented later in this manual.

Where does the name \emph{Twin-64} come from? The 64'' refers to the processor's native 64-bit architecture. The meaning of Twin’’ becomes apparent in the implementation chapters, where one possible realization of the architecture employs a dual-issue pipeline capable of executing two compatible instructions in parallel. This dual-issue organization represents the architectural idea behind the name, although the instruction set itself is not tied to a particular implementation. Simpler implementations, such as a single-issue processor with one arithmetic logic unit, are equally valid realizations of the Twin-64 architecture.


\unnumberedSection{Design Goals}

While it is not a design goal to build an industry-strength, competitive CPU, the CPU should nevertheless be useful and largely follow established design practices. For example, instructions should not be invented just because they seem useful, but rather designed with the assembler and compilers in mind. Instruction set design is CPU design. Twin-64 instructions will be hardwired and are in general pipeline-friendly, avoiding data and control hazards and stalls where possible.

\unnumberedSection{This book}

This document serves as a comprehensive guide to the Twin-64 processor, its instruction set, and runtime environment. It is organized into several parts, each focusing on a key aspect of the system.

\textbf{Part Two} introduces the Twin-64 architecture, providing an overview of the overall system and processor’s structure, core components, and design philosophy. This sets the stage for understanding how instructions are executed and how data moves through the system.

\textbf{Part Three} presents the instruction set in detail. It covers computational, immediate, memory reference, branch, and system control instructions. These chapters serve as the definitive reference for programmers and designers alike, offering precise information on encoding, operation, and usage.

\textbf{Part Four} discusses the runtime environment and software conventions. It explains program interaction with Twin-64, including calling conventions, exception handling, and runtime support for higher-level languages.

\textbf{Part Five} explores the I/O subsystem, detailing the mechanisms by which the processor communicates with external devices and handles input/output operations efficiently. ( ??? rather module as the overarching name ? )

\textbf{Part Six} presents the T64 firmware architecture.  ( ??? PAL, SAL, etc. Reset sequence, etc. )

\textbf{Part Seven} presents a reference implementation of the architecture and instruction set processor in a simulator environment.

Finally, the \textbf{Appendix} provides supplementary material, including additional design notes, tables, and clarifications to support the main text and offer deeper insight into the Twin-64 system.